% §8 Examples and moments — statements verbatim from blueprint part 08, which promotes
% the draft's Examples 8.1–8.3 to propositions (they are load-bearing: the locality
% theorem classifies the Bessel family, the ladder builds on the Gamma family) and
% splits the per-family clauses by cost. The paper keeps every node but REORDERS for
% reading: the admissibility criterion pair first (the one obligation every example
% shares, discharged once), then the moment formulas, then the three families with
% their split halves adjacent, then the heavy-tails remark. Adapted refs, flagged
% inline: prop:bessel-family's closing cross-reference (draft theorem numbers → the
% paper's sections), prop:gamma-kernels' proof (causal uniqueness → Prop. laplace-uniqueness).

\section{Examples and moments}\label{sec:examples}

Throughout this section we work in the canonical gauge, so that
$\phi_x = \Lap^{-1}[e^{-F(x s)}]$, and $T_x$ denotes a random variable with density $\phi_x$.
Each family below is admissible by one criterion, stated once; its converse makes the two
integrability conditions characteristic of the class rather than an artifact of the estimate.

% shared with blueprint prop:admissibility-criterion
\begin{proposition}[admissibility criterion for a memory kernel]\label{prop:admissibility-criterion}
Let $b_0 \ge 0$ and let $k : (0,\infty) \to [0,\infty)$ be nonincreasing with
\[
  \int_0^1 k(t)\,dt < \infty
  \qquad\text{and}\qquad
  \int_1^\infty \frac{k(t)}{t}\,dt < \infty .
\]
Then $F(s) = b_0 s + \int_0^\infty (1 - e^{-st})\,k(t)\,\frac{dt}{t}$ is finite for every
$s \ge 0$.

Consequently $F$ is of the form (7.1) and Theorem~\ref{thm:main-characterization} applies to it
--- a consequence and not part of the statement, which is why this node's dependency list records
only Definition~\ref{def:levy-exponent}: convergence is all the estimates below establish, and
admissibility is what Lemma~\ref{lem:selfdecomposable-exponents} then reads off it.
\end{proposition}

\begin{proof}
Split the integral at $t = 1$. On $(0,1]$ use $1 - e^{-st} \le st$, which is $1 + x \le e^x$ at
$x = -st$; the integrand is then at most $s\,k(t)$, integrable by the first hypothesis. On
$(1,\infty)$ use $1 - e^{-st} \le 1$; the integrand is at most $k(t)/t$, integrable by the
second. The drift term is finite for every real $s$.
\end{proof}

% shared with blueprint lem:criterion-converse
\begin{lemma}[the criterion is characteristic]\label{lem:criterion-converse}
Conversely, if $F$ is of the form (7.1) with $F(1) < \infty$, then $\int_0^1 k(t)\,dt < \infty$
and $\int_1^\infty k(t)t^{-1}dt < \infty$.
\end{lemma}

\begin{proof}
Take $s = 1$ in (7.1), so that $\int_0^\infty (1 - e^{-t})k(t)t^{-1}dt < \infty$. For $t > 0$,
$1 - e^{-t} \ge t/(1+t)$, so the integrand is at least $k(t)/(1+t)$. On $(0,1]$ that is at least
$k(t)/2$, and on $(1,\infty)$ it is at least $k(t)/(2t)$; both claims follow.
\end{proof}

In particular every admissible $F$ carries both integrability facts already, so nothing below
assumes regularity of $k$ beyond the definition; this is also what justified differentiating
under the integral sign in Lemma~\ref{lem:memory-kernel}.

% shared with blueprint prop:moments
\begin{proposition}[the mean delay and the influence curve]\label{prop:moments}
$\Ex\,T_x = x\,F'(\zp) = x\bigl(b_0 + \int_0^\infty k(t)\,dt\bigr)$, finite if and only if $k$
is integrable at infinity. When the mean is finite, the influence curve of
\sscite{fagerstrom2005temporal}, \S7, may be defined by the mean delay, and it is exactly
linear in the canonical gauge:
\[
  t_0 - t_m = F'(\zp)\; x .
\]
\end{proposition}

\begin{proof}
Differentiating the transform at the origin,
$\Ex\,T_x = -\partial_s e^{-F(xs)}\big|_{s=\zp} = x F'(\zp)$, and monotone convergence in
(7.1) gives $F'(\zp) = b_0 + \int_0^\infty k(t)\,dt$, finite exactly when $k$ is integrable at
infinity. The influence curve is then linear in $x$ because $\Ex\,T_x$ is.
\end{proof}

The higher moments run through one cited equivalence, the moment criterion for infinitely
divisible laws.

% shared with blueprint prop:moment-criterion
\begin{proposition}[the moment criterion]\label{prop:moment-criterion}
$\Ex\,T_x^{\,n} < \infty$ if and only if $\int_1^\infty t^{\,n-1} k(t)\,dt < \infty$.
\end{proposition}

\begin{proof}
The Lévy measure of $T_x$ is $\nu(dt) = k(t)t^{-1}dt$ up to the scaling by $x$, so
$\int_{t>1} t^n\,\nu(dt) = \int_1^\infty t^{\,n-1}k(t)\,dt$; the moment criterion for
infinitely divisible laws (\sscite{sato1999levy}, Thm.~25.3) equates finiteness of that
integral with finiteness of $\Ex\,T_x^n$.
\end{proof}

% shared with blueprint prop:stable-family
\begin{proposition}[the extremal stable family: the exponent]\label{prop:stable-family}
For $0 < \alpha < 1$ the exponent $F(s) = s^\alpha$ is of the form (7.1), with Lévy density
factor
\[
  k(t) = \frac{\alpha}{\Gamma(1-\alpha)}\, t^{-\alpha},
\]
which is nonincreasing.
\end{proposition}

\begin{proof}
$s^\alpha = \frac{\alpha}{\Gamma(1-\alpha)}\int_0^\infty (1 - e^{-st})\,t^{-\alpha-1}\,dt$ is
the standard integral representation, which is (7.1) with $b_0 = 0$ and
$k(t)/t = \alpha t^{-\alpha-1}/\Gamma(1-\alpha)$; $k$ is a negative power, hence
nonincreasing.
\end{proof}

% shared with blueprint prop:stable-moments
\begin{proposition}[the extremal stable family: every moment is infinite]\label{prop:stable-moments}
With $F$ as in Proposition~\ref{prop:stable-family}, $\Ex\,T_x^{\,n} = \infty$ for every $n \ge 1$
and every $x > 0$, so the delay of the 2005 stable kernels cannot be measured by any moment.
\end{proposition}

\begin{proof}
By Proposition~\ref{prop:moments} the mean delay is $\Ex\,T_x = x\bigl(b_0 + \int_0^\infty
k\bigr)$, finite exactly when $k$ is integrable at infinity; here $b_0 = 0$ and
$k(t) = \alpha t^{-\alpha}/\Gamma(1 - \alpha)$, which is not integrable at infinity for
$\alpha < 1$. So $\Ex\,T_x = \infty$. Every higher moment follows without a further criterion:
$t \le 1 + t^{\,n}$ for $t \ge 0$ and $n \ge 1$, so on the probability measure $\mu_{0,x}$ we get
$\Ex\,T_x \le 1 + \Ex\,T_x^{\,n}$, and a divergent mean forces $\Ex\,T_x^{\,n} = \infty$.
\end{proof}

% shared with blueprint prop:stable-mode
\begin{proposition}[the extremal stable family: the mode]\label{prop:stable-mode}
With $F$ as in Proposition~\ref{prop:stable-family}, the mode of $T_x$ scales as $x^{1/\alpha}$.
\end{proposition}

\begin{proof}
The self-similarity is the case $\chi(x) = x^{1/\alpha}$ of Corollary~\ref{cor:semigroup-case},
and the mode of $x^{1/\alpha}T_1$ is $x^{1/\alpha}$ times the mode of $T_1$.
\end{proof}

% shared with blueprint prop:gamma-family
\begin{proposition}[the Gamma family: the exponent]\label{prop:gamma-family}
For $\gamma > 0$ the exponent $F(s) = \gamma\,\log(1 + s)$ is of the form (7.1) with
$k(t) = \gamma\, e^{-t}$, which is nonincreasing.
\end{proposition}

\begin{proof}
The Frullani integral gives
$\log(1+s) = \int_0^\infty (1 - e^{-st})\,e^{-t}\,t^{-1}\,dt$, so $F$ has the form (7.1) with
$b_0 = 0$ and $k(t) = \gamma e^{-t}$, nonincreasing; hence $F$ is admissible by
Lemma~\ref{lem:selfdecomposable-exponents}.
\end{proof}

% shared with blueprint prop:gamma-kernels
\begin{proposition}[the Gamma family: the transforms and the mean]\label{prop:gamma-kernels}
With $F$ as in Proposition~\ref{prop:gamma-family}, the kernels and the increment measures have
transforms
\[
  e^{-F(xs)} = (1 + x s)^{-\gamma}, \qquad
  e^{-g_{a,b}(s)} = \Bigl(\tfrac{1 + as}{1 + bs}\Bigr)^{\gamma},
\]
and the mean delay is $\Ex\,T_x = \gamma x$.
\end{proposition}

\begin{proof}
$F(s) = \gamma\log(1+s)$ by Proposition~\ref{prop:gamma-family}, so $e^{-F(xs)} =
(1+xs)^{-\gamma}$; and $g_{a,b} = F(b\cdot) - F(a\cdot)$, so the increment transform is the ratio
$e^{-F(bs)}/e^{-F(as)}$, which is the stated one. For the mean, Proposition~\ref{prop:moments}
gives $\Ex\,T_x = x\bigl(b_0 + \int_0^\infty k\bigr)$ with $b_0 = 0$ and $k(t) = \gamma e^{-t}$,
whose integral is $\gamma$.
\end{proof}

% shared with blueprint prop:gamma-density
\begin{proposition}[the Gamma family: the kernel density]\label{prop:gamma-density}
With $F$ as in Proposition~\ref{prop:gamma-family} and $\gamma > 0$,
\[
  \phi_x(t) = \frac{t^{\gamma - 1} e^{-t/x}}{\Gamma(\gamma)\, x^{\gamma}},
\]
the Gamma density of fixed shape $\gamma$ and scale $x$.
\end{proposition}

\begin{proof}
$(1+xs)^{-\gamma}$ is the Laplace transform of the stated Gamma density --- the Gamma integral
$\int_0^\infty t^{\gamma-1}e^{-(1/x + s)t}\,dt = \Gamma(\gamma)/(1/x + s)^{\gamma}$ --- so
Proposition~\ref{prop:laplace-uniqueness} identifies $\phi_x$.
\end{proof}

% Verbatim from the blueprint's rem:gamma-sato; not shared, since remarks are not shareable.
\begin{remark}[the Gamma family as a Sato subordinator]\label{rem:gamma-sato}
The increment measures $\mu_{a,b}$ of this family are the increments of the Gamma-type Sato
subordinator. This is Remark~\ref{rem:sato-reading} read at the Gamma exponent.
\end{remark}

% shared with blueprint prop:gamma-moments
\begin{proposition}[the Gamma family: the higher moments]\label{prop:gamma-moments}
With $F$ as in Proposition~\ref{prop:gamma-family}, every moment of $T_x$ is finite, and
$\operatorname{Var} T_x = \gamma x^2$.
\end{proposition}

\begin{proof}
By Proposition~\ref{prop:gamma-density} the law of $T_x$ is the Gamma law of shape $\gamma$ and
rate $1/x$, so $\Ex\,T_x^{\,n} = \Gamma(\gamma+n)x^{n}/\Gamma(\gamma)$, finite for every $n$. In
particular $\Ex\,T_x = \gamma x$ and $\Ex\,T_x^{2} = \gamma(\gamma+1)x^2$, so
$\operatorname{Var} T_x = \gamma(\gamma+1)x^2 - \gamma^2x^2 = \gamma x^2$.
\end{proof}

% shared with blueprint prop:bessel-family — closing cross-reference now a \ref to the
% node the blueprint's hard-coded theorem numbers (10.4, 11.6, 12.5) meant.
\begin{proposition}[the Bessel-K family; preview of the memory line]\label{prop:bessel-family}
For $a > 0$ let $T_1 \stackrel{d}{=} 1/(2\gamma_a)$ be the scaled inverse-gamma law, where
$\gamma_a$ has the Gamma density of shape $a$, and let $F$ be its Laplace exponent. Then
\[
  e^{-F(s)} = \frac{2^{1-a}}{\Gamma(a)}\,\bigl(\sqrt{2s}\bigr)^{a} K_a\bigl(\sqrt{2s}\bigr),
\]
and $T_1$ is self-decomposable, so $F$ is of the form (7.1) and the family is admissible. In
the parabolic gauge $\tilde x = x^2/2$ these are the boundary-hitting kernels of the
Bessel-type media $\tfrac12\partial_x^2 + \tfrac{\beta}{x}\partial_x$ with $a = \tfrac12 -
\beta$, and the moments satisfy $\Ex\,T_x^n < \infty$ if and only if $a > n$. This family is
the subject of the memory-line development, \ref{thm:scale-cauchy}, \ref{thm:signaling-form} and
\ref{thm:locality}.
\end{proposition}

\begin{proof}
Write $\gamma_a$ for a Gamma$(a,1)$ variable and $T_1 = 1/(2\gamma_a)$. Substituting
$u = 1/(2t)$ in the Gamma density gives the density of $T_1$,
\[
  \phi_1(t) = \frac{t^{-a-1}\,e^{-1/(2t)}}{2^{a}\,\Gamma(a)}, \qquad t > 0 .
\]
The classical integral representation
$\int_0^\infty t^{\nu-1}\exp(-A/t - Bt)\,dt = 2\,(A/B)^{\nu/2}K_\nu\bigl(2\sqrt{AB}\bigr)$,
taken at $\nu = -a$, $A = \tfrac12$, $B = s$, together with $K_{-a} = K_a$, gives
\[
  \Ex\,e^{-sT_1}
  = \frac{1}{2^{a}\Gamma(a)}\int_0^\infty t^{-a-1}e^{-st - 1/(2t)}\,dt
  = \frac{2\,(2s)^{a/2}\,K_a\bigl(\sqrt{2s}\bigr)}{2^{a}\Gamma(a)}
  = \frac{2^{1-a}}{\Gamma(a)}\bigl(\sqrt{2s}\bigr)^{a}K_a\bigl(\sqrt{2s}\bigr).
\]
Self-decomposability of $T_1$ is a theorem of Halgreen (\sscite{halgreen1979self}, for the
generalized inverse Gaussian laws, of which this is a limiting case); by
Lemma~\ref{lem:selfdecomposable-exponents} its exponent is therefore of the form (7.1), and
Theorem~\ref{thm:main-characterization} makes the family admissible. Finally, the Lévy density
factor of this $F$ satisfies $k(t) \sim a$ as $t \to 0$ and decays like $t^{-a}$ at infinity,
so $\int_1^\infty t^{n-1}k(t)\,dt < \infty$ exactly when $n < a$;
Proposition~\ref{prop:moment-criterion} converts that into $\Ex\,T_x^n < \infty$ iff $a > n$.
\end{proof}

% Verbatim from the blueprint's rem:heavy-tails.
\begin{remark}[heavy tails are the semigroup's doing]\label{rem:heavy-tails}
The heavy tails of the 2005 kernels are an artifact of the semigroup axiom, not of causality
together with covariance. Within Theorem~\ref{thm:main-characterization} the semigroup sub-case
(Corollary~\ref{cor:semigroup-case}) is precisely the subset with $k$ a pure power --- the only
$k$ compatible with homogeneity --- and every non-power nonincreasing $k$, for instance
$k(t) = \gamma e^{-t}$, yields a causal, continuously scale-covariant family with finite
moments of all orders.
\end{remark}
